% 06-basics.tex — Clause 6: Basic concepts
\chapter{Basic concepts}
\label{sec:06-basics}
\index{basic concepts}

\section{Declarations and definitions}
\label{sec:basics.declarations}
\index{declaration}
\index{definition}

\pnum
A \textit{declaration} introduces a name into a scope and associates it with
an entity: a variable, function, procedure, type, or module.  Every name
shall be declared before it is used, with the exception of names that are
resolved by Universal Function Call Syntax (\S\,\ref{sec:12-functions}).

\pnum
A \textit{definition} is a declaration that also allocates storage (for
variables) or provides a body (for functions and procedures).  In Lain,
every declaration at the local level is also a definition.

\begin{constraint}
  Redeclaring a name already visible in the current or an enclosing scope
  --- a duplicate parameter name, or a local that repeats or shadows an
  outer binding --- is \illformed{}~(\diagcode{E013}).  Lain forbids
  shadowing (\S\,\ref{sec:basics.scopes}).
\end{constraint}

\begin{constraint}
  Use of an undeclared identifier is \illformed{}~(\diagcode{E106}).  The
  reference implementation issues this diagnostic when the translation output is
  generated --- after name resolution, monomorphization, and UFCS desugaring are
  complete --- so that a name still unbound at that point is unambiguously
  undeclared.
\end{constraint}

\begin{example}
\begin{laincode}
proc f() {
    var x = 1
    var x = 2   // [E013]: redeclaration of 'x' is forbidden
}
\end{laincode}
\end{example}

\section{Scopes}
\label{sec:basics.scopes}
\index{scope}

\pnum
A \textit{scope} is a region of program text within which a declared name
can be used.  Lain has the following scope levels:

\begin{description}[leftmargin=2em]
  \item[Module scope]
    Names declared at the top level of a source file are in module scope.
    Module-scope names are accessible throughout the module and, by
    default, accessible from other modules that import this module.

  \item[Function/procedure scope]
    Parameter names are in the scope of the function or procedure body.

  \item[Block scope]
    Names declared inside a block (\lain{\{ \}}) are in the scope of that
    block.  They are not accessible outside the enclosing block.
\end{description}

\pnum
Scopes are \textit{lexically nested}.  Lain \emph{forbids shadowing}
entirely: a declaration whose name is already visible in the current scope
or in any enclosing scope --- a function parameter shadowed by a local, or a
name redeclared in a nested block --- is \illformed{}~(\diagcode{E013}).

\begin{rationale}
  Disallowing shadowing removes a class of bugs in which an inner binding
  silently captures a name the reader expects to denote an outer one.  Every
  identifier in a function body therefore resolves to exactly one
  declaration.
\end{rationale}

\pnum
A conforming implementation shall support at least 64 levels of nested
block scope (\S\,\ref{sec:conformance.limits}).

\section{Name resolution}
\label{sec:basics.resolution}
\index{name resolution}

\pnum
Name resolution maps an identifier in source text to its declaration.
The lookup algorithm is:

\begin{enumerate}
  \item Search the innermost block scope.
  \item If not found, search progressively wider enclosing block scopes.
  \item If not found, search the function/procedure scope (parameters).
  \item If not found, search the module scope.
  \item If not found, search imported module scopes (in import order).
  \item If not found, attempt UFCS resolution (\S\,\ref{sec:12-functions}).
  \item If still not found, the identifier is undeclared and the program is
        \illformed{}~(\diagcode{E106}).
\end{enumerate}

\pnum
Because shadowing is forbidden (\S\,\ref{sec:basics.scopes}), a name resolves
to at most one declaration along the scope chain: a lookup that would find a
second binding of the same name is rejected at the point of the redeclaration
(\diagcode{E013}), not silently resolved by an innermost-wins rule.

\section{Storage duration}
\label{sec:basics.storage}
\index{storage duration}

\pnum
Every object has a \textit{storage duration} that determines the lifetime
of its storage.

\begin{description}[leftmargin=2em]
  \item[Automatic storage duration]
    Variables declared inside a block or function body.  Storage is
    allocated on entry to the block and deallocated on exit.

  \item[Static storage duration]
    Variables declared at module scope.  Storage persists for the lifetime
    of the program.  \idbehav{whether module-scope variables are zero-
    initialized before \lain{proc main} is called}.

  \item[Arena storage duration]
    Objects allocated from an explicit \lain{Arena}.  Storage persists
    until the arena is freed as a whole.
\end{description}

\section{Objects and values}
\label{sec:basics.objects}
\index{object}

\pnum
An \textit{object} is a region of storage interpreted according to a type.
Every declared variable corresponds to an object.

\pnum
Every variable has:
\begin{itemize}
  \item a \textit{declared type} (\S\,\ref{sec:07-types});
  \item an \textit{ownership mode}: owned (default for linear types),
        \lain{shared} (read-only reference), or \lain{var} (mutable
        reference) (\S\,\ref{sec:11-ownership});
  \item a \textit{linear state}: Uninit, Unconsumed, or Consumed (for
        linear types only) (\S\,\ref{sec:11-ownership}).
\end{itemize}

\section{Lvalues and rvalues}
\label{sec:basics.lvalues}
\index{lvalue}
\index{rvalue}

\pnum
An \textit{lvalue} is an expression that designates an object with a
persistent storage location: a variable identifier, a field access on an
lvalue, or an array index expression on an lvalue.

\pnum
An \textit{rvalue} is an expression whose value has no persistent storage
location: literals, the result of arithmetic operations, or the result of
a function call.

\pnum
Assignment requires the left-hand side to be an lvalue.

\begin{constraint}
  An assignment whose left-hand side is an rvalue is
  \illformed{}~(\diagcode{E100}).
\end{constraint}

\section{Alignment}
\label{sec:basics.alignment}
\index{alignment}

\pnum
Every type has a required alignment.  The alignment of a type is
\idbehav{the same as the alignment of the corresponding C99 type}.
Struct layout follows C99 struct layout rules, including padding
(\S\,\ref{sec:07-types}).

\pnum
Stack-allocated variables are aligned to their type's required alignment.
The implementation shall not insert additional padding between local
variables unless required by the hardware.
